% Example "robertson": the problem, the concentrations and the fast eigenvalue,
% and the stability tests of explicit methods on a stiff problem. Every number
% comes from data/, written by examples/robertson/main.f90: run
% `fpm run --example robertson` from the project root, then build this file
% with lualatex in examples/robertson/.
\documentclass[11pt,a4paper]{article}
\usepackage[margin=2.2cm]{geometry}
\usepackage{amsmath,booktabs,array}
\usepackage{pgfplots}
\usepgfplotslibrary{groupplots}
\pgfplotsset{compat=1.18}
\pagestyle{empty}

\colorlet{m-euler}{black}
\colorlet{m-order2}{orange!85!black}
\colorlet{m-order3}{green!50!black}
\colorlet{m-rk4}{blue!80!black}
\colorlet{m-dp54}{red!75!black}

\begin{document}
\section*{Example: Robertson's reaction, a stiff problem}

\subsection*{1. Problem}
From the SciML problem library (\texttt{prob\_ode\_rober}), Robertson's
model of three chemical reactions with rates of very different sizes:
\begin{equation*}
  y_1' = -0.04\,y_1 + 10^4 y_2 y_3, \qquad
  y_2' = 0.04\,y_1 - 10^4 y_2 y_3 - 3\cdot10^7 y_2^2, \qquad
  y_3' = 3\cdot10^7 y_2^2, \qquad y(0) = (1, 0, 0).
\end{equation*}
The total $y_1 + y_2 + y_3 = 1$ is conserved. $y_2$ rises within
$10^{-3}$ time units to a tiny quasi-steady value, after which everything
changes slowly: $y_1$ falls to half its value only at $t \approx 270$.
The Jacobian has the eigenvalue $0$ (the conservation) and two others,
$\lambda = \bigl(T \pm \sqrt{T^2 - 4D}\bigr)/2$ with $T$ its trace and $D$
the sum of its principal $2\times2$ minors; the fast one, $\lambda_f$, reaches
$-2.2\times10^3$ within $4\times10^{-3}$ time units and stays between that and
$-6.8\times10^3$ up to $t = 1000$. An explicit step is
stable only while $|\lambda_f(t)|\,\Delta t \le h^*$: it must stay below about
$10^{-3}$ for as long as the integration lasts, although nothing in the
solution moves on that time scale. That is stiffness.
\begin{center}
\small
% Written by examples/robertson/main.f90: do not edit.
\begin{tabular}{l>{\raggedright\arraybackslash}p{12.4cm}}
\toprule
rate constants & $k_1 = 0.04$, $k_2 = 3\times10^7$, $k_3 = 10^4$ \\
initial condition & $y = (1, 0, 0)$; interval $0 \le t \le 1000$ \\
reference & \texttt{dp54}, $\Delta t = 5\times10^{-6}$; 400 samples, log-spaced from $t = 10^{-6}$ \\
fast eigenvalue & $|\lambda_f| = |T - \sqrt{T^2 - 4D}|/2$, $T = -0.04 - 10^4 y_3 - 6\times10^7 y_2$, $D = 2.4\times10^6 y_2 + 6\times10^{11} y_2^2$ \\
step sweep & 14 steps, log-spaced from $10^{-4}$ to $10^{-2}$ \\
failure & first $t$ with a value not finite or outside $[-0.01, 1.01]$, $y_2 < -10^{-6}$, or (for $t > 10^{-3}$) $y_2$ off the reference by more than half \\
predicted onset & first sampled $t$ with $|\lambda_f(t)|\,\Delta t > h^*$ \\
reaching $T$ & largest $\Delta t$ whose run reaches $T$ without blowing up; bisection (14 halvings) in $[0.7, 1.5]$ times $h^*/\max_{t \le T}|\lambda_f|$ \\
figures & \texttt{rk4} at $\Delta t = 2\times10^{-3}$, $7\times10^{-4}$, $3\times10^{-4}$ \\
methods & \texttt{rk3ws}, \texttt{rk4}, \texttt{dp54} \\
interface & plain arrays, \texttt{rk\_array\_scheme\_t} \\
\bottomrule
\end{tabular}

\end{center}

\newpage
\subsection*{2. Solutions}
\begin{center}
\begin{tikzpicture}
  \begin{groupplot}[
      group style={group size=1 by 2, vertical sep=1.6cm}, % two panels, one above the other
      width=14cm, height=7.2cm,                            % wide: time spans nine decades
      xmode=log, xmin=1e-6, xmax=1000,                     % logarithmic time
      grid=major, grid style={black!8},                    % faint grid
      axis background/.style={fill=black!2},               % very light background
    ]
    % ---- (a) concentrations ----
    \nextgroupplot[
      title={(a) concentrations, and \texttt{rk4} at three fixed steps}, ymin=0, ymax=1.05,
      ylabel={concentration}, xticklabels={},              % shares the t axis with (b)
      legend style={font=\small, at={(0.02,0.5)}, anchor=west, fill=white, fill opacity=0.85, text opacity=1, draw=black!30},
    ]
    \addplot[black, very thick] table[x=t, y=y1] {data/reference.dat};                       % y1
    \addlegendentry{$y_1$}
    \addplot[black!50, very thick, dashed] table[x=t, y expr=1e4*\thisrow{y2}] {data/reference.dat}; % y2 x 1e4
    \addlegendentry{$10^4\,y_2$}
    \addplot[black, very thick, dotted] table[x=t, y=y3] {data/reference.dat};               % y3
    \addlegendentry{$y_3$}
    \addplot[red!80!black, thick, mark=*, mark size=1pt] table[x=t, y expr=1e4*\thisrow{y2}] {data/fixed_1.dat}; % rk4, 2e-3
    \addlegendentry{\texttt{rk4}, $\Delta t = 2\cdot10^{-3}$ ($10^4 y_2$)}
    \addplot[orange!85!black, thick] table[x=t, y expr=1e4*\thisrow{y2}] {data/fixed_2.dat};  % rk4, 7e-4
    \addlegendentry{\texttt{rk4}, $\Delta t = 7\cdot10^{-4}$}
    \addplot[m-rk4, thick] table[x=t, y expr=1e4*\thisrow{y2}] {data/fixed_3.dat};            % rk4, 3e-4
    \addlegendentry{\texttt{rk4}, $\Delta t = 3\cdot10^{-4}$}
    % ---- (b) the fast eigenvalue and the thresholds ----
    \nextgroupplot[
      title={(b) $|\lambda_f(t)|$ and the limits $h^*/\Delta t$ of the three \texttt{rk4} runs},
      ymode=log, ymin=1, ymax=2e4, xlabel={$t$}, ylabel={$|\lambda_f|$},
    ]
    \addplot[black, very thick] table[x=t, y=lambda] {data/reference.dat};                  % fast eigenvalue
    \addplot[red!80!black, thick, dashed, domain=1e-6:1000] {2.7853/2e-3};                  % rk4 limit, dt = 2e-3
    \addplot[orange!85!black, thick, dashed, domain=1e-6:1000] {2.7853/7e-4};               % rk4 limit, dt = 7e-4
    \addplot[m-rk4, thick, dashed, domain=1e-6:1000] {2.7853/3e-4};                         % rk4 limit, dt = 3e-4
    \node[font=\scriptsize, text=red!80!black, anchor=south west] at (axis cs:1.5e-6,1392) {$h^*/(2\cdot10^{-3})$};
    \node[font=\scriptsize, text=orange!85!black, anchor=south west] at (axis cs:1.5e-6,3979) {$h^*/(7\cdot10^{-4})$};
    \node[font=\scriptsize, text=blue!80!black, anchor=south west] at (axis cs:1.5e-6,9284) {$h^*/(3\cdot10^{-4})$};
  \end{groupplot}
\end{tikzpicture}
\end{center}
(a) The classic picture on a logarithmic time axis: $y_1$ decays into $y_3$
through the short-lived intermediate $y_2$ (scaled by $10^4$). \texttt{rk4}
with $\Delta t = 2\cdot10^{-3}$ fails after seven steps (red, $t = 0.014$); with
$7\cdot10^{-4}$ it follows the reference until $t \approx 99$ and then fails;
only with $3\cdot10^{-4}$ does it reach $t \approx 960$. (b) Why: a run is
stable while the fast eigenvalue (black) stays below its limit $h^*/\Delta t$
(dashed, $h^* = 2.785$ for \texttt{rk4}). $|\lambda_f|$ jumps to about
$2190$ within $4\cdot10^{-3}$ time units and then creeps up as $y_3$ grows;
each run fails shortly after the black curve crosses its line.

\newpage
\subsection*{3. Tests: where the explicit step fails}

\paragraph{Prediction.} A run with step $\Delta t$ can follow the solution
only while $|\lambda_f(t)|\,\Delta t \le h^*$. (i) For a given step it fails
soon after the predicted onset, the first time this is violated. (ii) The
largest step that reaches a time $T$ is $h^*/\max_{t \le T}|\lambda_f(t)|$.
(iii) To reach $t = 1000$ an explicit method needs about
$1000\,\max|\lambda_f|/h^*$ steps, although the solution changes on a time
scale of tens or hundreds of units.

\paragraph{Result.}
\begin{center}
\small
% Written by examples/robertson/main.f90: do not edit.
\begin{tabular}{lcccccc}
\toprule
& & \multicolumn{4}{c}{largest step reaching $T$, measured / predicted} & \\
\cmidrule(lr){3-6}
name & $h^*$ & $T = 1$ & $T = 10$ & $T = 100$ & $T = 1000$ & steps to $t = 1000$ \\
\midrule
\texttt{rk3ws} & 2.5127 & 1.0876 & 1.0172 & 0.9967 & 0.9921 & $2.69\times10^{6}$ \\
\texttt{rk4} & 2.7853 & 1.0590 & 1.1554 & 1.0104 & 0.9938 & $2.42\times10^{6}$ \\
\texttt{dp54} & 3.3066 & 0.8324 & 0.9604 & 1.0385 & 1.0722 & $2.04\times10^{6}$ \\
\bottomrule
\end{tabular}

\end{center}

\begin{center}
\begin{tikzpicture}
  \begin{groupplot}[
      group style={group size=2 by 1, horizontal sep=1.8cm}, % two panels
      width=7.4cm, height=6.8cm,                           % size of each panel
      xmode=log, ymode=log, grid=major, grid style={black!8}, % log-log
      axis background/.style={fill=black!2},               % very light background
    ]
    % ---- (i) failure time against the step ----
    \nextgroupplot[
      title={(i) failure time: measured (marks), onset (lines)}, xlabel={$\Delta t$}, ylabel={$t$},
      legend to name=roblegend, legend columns=5,          % legend below the panels
      legend style={font=\small, /tikz/every even column/.append style={column sep=0.3cm}},
    ]
    \addplot[m-rk4, thick] table[x=dt, y=t] {data/onset_rk4.dat};
    \addlegendentry{\texttt{rk4}}
    \addplot[m-rk4, only marks, mark=diamond*, mark size=1.9pt, forget plot] table[x=dt, y=t] {data/blowup_rk4.dat};
    \addplot[m-dp54, thick] table[x=dt, y=t] {data/onset_dp54.dat};
    \addlegendentry{\texttt{dp54}}
    \addplot[m-dp54, only marks, mark=pentagon*, mark size=1.9pt, forget plot] table[x=dt, y=t] {data/blowup_dp54.dat};
    % ---- (ii) largest step reaching T ----
    \nextgroupplot[
      title={(ii) largest $\Delta t$ reaching $T$}, xlabel={$T$}, ylabel={$\Delta t$},
    ]
    \addplot[m-rk4, thick] table[x=T, y=predicted] {data/reach_rk4.dat};
    \addplot[m-rk4, only marks, mark=diamond*, mark size=1.9pt] table[x=T, y=measured] {data/reach_rk4.dat};
    \addplot[m-dp54, thick] table[x=T, y=predicted] {data/reach_dp54.dat};
    \addplot[m-dp54, only marks, mark=pentagon*, mark size=1.9pt] table[x=T, y=measured] {data/reach_dp54.dat};
  \end{groupplot}
  \node[below=0.2cm] at (current bounding box.south)        % legend centred under the panels
    {\pgfplotslegendfromname{roblegend}};
\end{tikzpicture}
\end{center}
(i) Above about $1.7\cdot10^{-3}$ every method fails at its first step;
below, the measured failures (marks) follow the predicted onset (lines):
often within a few per cent (\texttt{rk4} at $\Delta t = 5.9\cdot10^{-4}$:
onset $171.0$, failure $178.8$), later where $|\lambda_f|$ crosses the limit slowly
and the instability needs time to grow. (ii) The largest step that reaches
$T$ agrees with $h^*/\max_{t\le T}|\lambda_f|$ within about $1\,\%$ for
$T = 100$ and $1000$ (table), except \texttt{dp54}, $4$--$7\,\%$ above: its
$R(-h)$ passes $+1$, not $-1$, at its limit, so an instability grows without
changing sign and the stiff term $3\cdot10^7y_2^2$ holds $y_2$ at a wrong but
bounded value; the run is caught only once $y_2$ is off by half. For $T = 1$
an instability that starts near the end has no time to show, and the
measured step is a little larger. Reaching $t = 1000$ takes two to three and
a half million steps, whatever the method; an implicit method follows the slow solution
with steps of the size of the solution's own time scale.
\end{document}
